\section{Closed forms for the touching cells}
\label{sec:touching}

Within the note, this section combines three earlier results: the
two-dimensional form of a moment (Theorem~\theoRef{reduction}), the weights that are polynomials on $I_-$ and $I_+$
(Corollary~\coroRef{weights}), and the corner difference with the closed-form antiderivatives $\Gf{p}{m}{n}$
(Proposition~\propRef{gproperties} and Appendix~\ref{app:antiderivatives}). In the literature, closed forms of static
potential integrals are classical. Wilton et al.~\cite{wilton1984potential} give them for uniform and linear sources on
polygons, a single integration over the source domain, and Tzoulis and Eibert~\cite{Tzoulis2004} review this line of
work. Lopez-Pena and Mosig~\cite{lopezpena2009analytical} evaluate the full fourfold static integrals of uniform sources
over rectangular cells, and Bao et al.~\cite{Bao2020} and the thesis~\cite{machacek2025method} do so for rooftop
functions on the coincident cell. This section extends the last two to the linear weights of the rooftops on all nine
touching cells, the edge and vertex neighbours included, and from the kernel $R^{-1}$ to the kernels $R$ and $R^3$ of
the three-term extraction.

Theorem~\theoRef{reduction} turns a moment into the two-dimensional integral~\eqref{eq:reduction} over $[-1,1]^2$,
with the weights $W_{\alpha\beta}(\nu)$, $W_{00}(\eta)$ and the distance $\rho_{\offs}(\nu,\eta)$. For a touching offset
$\offs=(d_x,d_y)\in\{0,1\}^2$ the distance vanishes at the singular point $(\nu,\eta)=(d_x,d_y)$, and this is the only
point where the integrand is not smooth. The plan of this section is to cut $[-1,1]^2$ into pieces on which the weights
are polynomials and to move the singular point of every piece to the origin. Each piece then becomes a rectangle in the
quadrant $Q$ of Section~\ref{sec:antiderivatives}, with polynomial weights in front of $(u^2+v^2)^{p/2}$, and
Proposition~\propRef{gproperties}(ii) integrates it exactly with the antiderivatives $\Gf{p}{m}{n}$.

Figure~\ref{fig:rooftoptouch} shows where these configurations come from. A rooftop consists of two wedges, so a pair
of rooftops gives four wedge pairs (Proposition~\propRef{generators}), and each wedge pair is a pair of cells. Two
rooftops in a row that overlap by one cell have a wedge pair on the same cell, the coincident configuration. With one
cell between them, the two inner wedges share an edge parallel to $y$. Two rooftops side by side in $y$ have wedge
pairs that share an edge parallel to $x$, and two rooftops placed diagonally have a wedge pair that meets in a single
vertex. In all four cases the two cells of the pair have the points of the red set in common, where $R=0$.

\begin{figure}[ht]
    \centering
    % Two x-directed rooftops whose wedges form a touching cell pair, one panel per configuration of Section touching: the
% set where the two cells meet (a cell, an edge or a vertex) is drawn in red on the floor.
\tdplotsetmaincoords{62}{20}
\tikzset{
    touchpanel/.style={tdplot_main_coords, line cap=round, line join=round, scale=0.88},
    touchset/.style={PairedF, line width=2.2pt},
}
\newcommand{\touchfloor}[2]{%
    \foreach \i in {0,...,#1} { \draw[dashed] ({\i*\dx},0,0) -- ({\i*\dx},{#2*\dy},0); }
    \foreach \j in {0,...,#2} { \draw[dashed] (0,{\j*\dy},0) -- ({#1*\dx},{\j*\dy},0); }%
}
% \touchwedge{line}{fill}{x0}{z0}{x1}{z1}: the strip of an x-directed wedge between x0*\dx and x1*\dx, with the heights
% z0*\roofh and z1*\roofh there. In panel (a) the two wedges on the shared cell cross at its midline, so each is split
% there and the halves are drawn back to front.
\newcommand{\touchwedge}[6]{%
    \fill[#2] ({#3*\dx},0,{#4*\roofh}) -- ({#5*\dx},0,{#6*\roofh}) -- ({#5*\dx},\dy,{#6*\roofh})
        -- ({#3*\dx},\dy,{#4*\roofh}) -- cycle;
    \draw[#1, thick] ({#3*\dx},0,{#4*\roofh}) -- ({#5*\dx},0,{#6*\roofh}) -- ({#5*\dx},\dy,{#6*\roofh})
        -- ({#3*\dx},\dy,{#4*\roofh}) -- cycle;%
}
\def\dx{1.3}
\def\dy{1.3}
\def\roofh{1.1}
\begin{tabular}{@{}c@{\hspace{0.6cm}}c@{}}
\begin{tikzpicture}[touchpanel]
    \touchfloor{3}{1}
    \touchwedge{PairedF}{PairedE}{0}{0}{1}{1}
    \touchwedge{PairedF}{PairedE}{1.5}{0.5}{2}{0}
    \touchwedge{PairedB}{PairedA}{1}{0}{1.5}{0.5}
    \touchwedge{PairedF}{PairedE}{1}{1}{1.5}{0.5}
    \touchwedge{PairedB}{PairedA}{1.5}{0.5}{2}{1}
    \touchwedge{PairedB}{PairedA}{2}{1}{3}{0}
    \foreach \i in {1,2} {
        \draw[thick] ({\i*\dx},0,0) -- ({\i*\dx},\dy,0);
        \draw[thick, dotted] ({\i*\dx},0,\roofh) -- ({\i*\dx},0,0);
        \draw[thick, dotted] ({\i*\dx},\dy,\roofh) -- ({\i*\dx},\dy,0);
    }
    \draw[touchset] (\dx,0,0) -- (2*\dx,0,0) -- (2*\dx,\dy,0) -- (\dx,\dy,0) -- cycle;
    \node[PairedF, anchor=east] at (0,0,{0.4*\roofh}) {$\V{\Psi}^x_{\V{0}}$};
    \node[PairedB, anchor=west] at ({3*\dx},{\dy},{0.3*\roofh}) {$\V{\Psi}^x_{(1,0)}$};
\end{tikzpicture}
&
\begin{tikzpicture}[touchpanel]
    \touchfloor{4}{1}
    \noteroofx{0}{0}{PairedF}{PairedE}{PairedF}{PairedE}
    \noteroofx{2}{0}{PairedB}{PairedA}{PairedB}{PairedA}
    \draw[touchset] (2*\dx,0,0) -- (2*\dx,\dy,0);
    \node[PairedF, anchor=east] at (0,0,{0.4*\roofh}) {$\V{\Psi}^x_{\V{0}}$};
    \node[PairedB, anchor=west] at ({4*\dx},{\dy},{0.3*\roofh}) {$\V{\Psi}^x_{(2,0)}$};
\end{tikzpicture}
\\
(a) coincident cell, $\offs=(0,0)$ & (b) shared edge, $\offs=(1,0)$
\\[0.5cm]
\begin{tikzpicture}[touchpanel]
    \touchfloor{2}{2}
    \noteroofx{0}{1}{PairedB}{PairedA}{PairedB}{PairedA}
    \noteroofx{0}{0}{PairedF}{PairedE}{PairedF}{PairedE}
    \draw[touchset] (0,\dy,0) -- (2*\dx,\dy,0);
    \node[PairedF, anchor=east] at (0,0,{0.4*\roofh}) {$\V{\Psi}^x_{\V{0}}$};
    \node[PairedB, anchor=west] at ({2*\dx},{2*\dy},{0.3*\roofh}) {$\V{\Psi}^x_{(0,1)}$};
\end{tikzpicture}
&
\begin{tikzpicture}[touchpanel]
    \touchfloor{4}{2}
    \noteroofx{2}{1}{PairedB}{PairedA}{PairedB}{PairedA}
    \noteroofx{0}{0}{PairedF}{PairedE}{PairedF}{PairedE}
    \fill[PairedF] (2*\dx,\dy,0) circle (3pt);
    \node[PairedF, anchor=east] at (0,0,{0.4*\roofh}) {$\V{\Psi}^x_{\V{0}}$};
    \node[PairedB, anchor=west] at ({4*\dx},{2*\dy},{0.3*\roofh}) {$\V{\Psi}^x_{(2,1)}$};
\end{tikzpicture}
\\
(c) shared edge, $\offs=(0,1)$ & (d) shared vertex, $\offs=(1,1)$
\end{tabular}

    \caption{Two $x$-directed rooftops, $\V{\Psi}^x_{\V{0}}$ (red) and a second one (blue), placed so that two of their
    wedges lie on touching cells; the set that the two cells have in common is drawn in red on the floor. (a) Rooftop
    offset $(1,0)$: the falling wedge of $\V{\Psi}^x_{\V{0}}$ and the rising wedge of $\V{\Psi}^x_{(1,0)}$ lie on the
    same cell. (b) Rooftop offset $(2,0)$: the wedges on $\Om{(1,0)}$ and $\Om{(2,0)}$ share an edge parallel to $y$.
    (c) Rooftop offset $(0,1)$: the wedges side by side share edges parallel to $x$. (d) Rooftop offset $(2,1)$: the
    wedges on $\Om{(1,0)}$ and $\Om{(2,1)}$ share only a vertex.}
    \label{fig:rooftoptouch}
\end{figure}

The cut is the one of Corollary~\coroRef{weights}: in each direction the interval $[-1,1]$ splits into $I_-=[-1,0]$ and
$I_+=[0,1]$, and on each of them the weights are polynomials. This gives the four pieces $I_s\times I_t$,
$s,t\in\{-,+\}$, of Figure~\ref{fig:regions}.

Moving the singular point to the origin is what makes one table of closed forms enough. First, at the origin the
kernel is $(u^2+v^2)^{p/2}$ without any shift, so every piece of every configuration is integrated by the same functions
$\Gf{p}{m}{n}$ of Definition~\defRef{antiderivative}, which depend only on the end point $(a,b)$ of the rectangle. A
singular point at some other position $(u_0,v_0)$ would need a separate family of antiderivatives for every shift.
Second, on a piece with $u,v\ge0$ the absolute values $\ABS{x-x'}$ and $\ABS{y-y'}$ in the distance and the kinks of the
weights at $\nu=d_x$ and $\eta=d_y$ disappear, so the integrand is a sum of terms $u^mv^n(u^2+v^2)^{p/2}$, the integrand
of~\eqref{eq:G}. Third, the singular point is a corner of every piece that contains it, and $\Gf{p}{m}{n}$ vanishes on
both axes (Proposition~\propRef{gproperties}(i)). The corner sum~\eqref{eq:corner} therefore never evaluates anything at
$R=0$, and the singularity is integrated exactly, without a limit. Placing the singularity at a vertex is also the
starting point of the Duffy transformation~\cite{duffy1982}, which uses it to make the integrand smooth for quadrature
instead of integrating it in closed form. The next lemma moves the singular point to the origin, one direction at a
time.

\lemmaTitle{The polynomial pieces}{pieces} Let $d\in\{0,1\}$ be the offset in the $x$-direction, so that the singular
value of $\nu$ is $\nu=d$. On each interval $I_s$, $s\in\{-,+\}$, introduce the distance from the singular value,
measured in units of length,
\begin{equation}
    \label{eq:udist}
    u = a\,\ABS{\nu-d} .
\end{equation}
On $I_s$ the difference $\nu-d$ does not change sign, so $\ABS{\nu-d}=\sigma(\nu-d)$ with a fixed sign
$\sigma\in\{-1,+1\}$, and the substitution is linear with the inverse $\nu=d+\sigma u/a$. Explicitly,
\begin{equation}
    \label{eq:ucases}
    \left\{
    \begin{array}{l@{\quad}r@{}l@{\qquad}r@{}l}
        d=0\colon & u & {}=-a\nu\in[0,a] \text{ on } I_-,     & u & {}=a\nu\in[0,a] \text{ on } I_+,\\[2pt]
        d=1\colon & u & {}=a(1-\nu)\in[a,2a] \text{ on } I_-, & u & {}=a(1-\nu)\in[0,a] \text{ on } I_+ .
    \end{array}
    \right.
\end{equation}
In all four cases the singular value $\nu=d$ corresponds to $u=0$. Every weight $W_{\alpha\beta}$ of~\eqref{eq:weights},
restricted to $I_s$ and written in $u$, is a polynomial in $u$ of degree at most three; these polynomials are listed in
Table~\ref{tab:pieces}. The $y$-direction is the same with $\eta$, $b$, $d_y$ and $v=b\,\ABS{\eta-d_y}$.

\vspace{\baselineskip}
\proof For $d=0$ the difference $\nu-d=\nu$ lies in $[-1,0]$ on $I_-$ and in $[0,1]$ on $I_+$, so $\sigma=-1$ and
$\sigma=+1$, respectively, and $u$ runs over $[0,a]$ on both. For $d=1$ the difference $\nu-1$ lies in $[-2,-1]$ on $I_-$
and in $[-1,0]$ on $I_+$, so $\sigma=-1$ on both, and $u=a(1-\nu)$ runs over $[a,2a]$ and $[0,a]$. This
gives~\eqref{eq:ucases}. On $I_s$ every weight is one of the polynomials of~\eqref{eq:weights}, of degree at most three
in $\nu$, and the linear substitution $\nu=d+\sigma u/a$ keeps it a polynomial of the same degree in $u$. Expanding gives
Table~\ref{tab:pieces}. \qed

\begin{figure}[ht]
    \centering
    % Top: the reduced domain [-1,1]^2 of Theorem reduction with the pieces I_s x I_t and the singular point (d_x,d_y).
% Bottom: the pieces as rectangles of the quadrant Q in u = s a (nu - d_x), v = t b (eta - d_y) (Lemma pieces), a = b.
\colorlet{pieceMM}{PairedA}
\colorlet{piecePM}{PairedG}
\colorlet{pieceMP}{PairedC}
\colorlet{piecePP}{PairedI}
\newcommand{\notechip}[4]{\node[fill=#3, inner sep=1pt, minimum width=0.34cm, font=\tiny] at (#1,#2) {$#4$};}
\begin{tikzpicture}[line cap=round, line join=round, scale=1.0]
    \foreach \col/\dxo/\dyo/\name in {0/0/0/{coincident, $\offs=(0,0)$}, 1/1/0/{edge, $\offs=(1,0)$},
                                     2/0/1/{edge, $\offs=(0,1)$}, 3/1/1/{vertex, $\offs=(1,1)$}} {
        \begin{scope}[shift={(3.2*\col,0)}]
            \fill[pieceMM] (-1,-1) rectangle (0,0);
            \fill[piecePM] (0,-1) rectangle (1,0);
            \fill[pieceMP] (-1,0) rectangle (0,1);
            \fill[piecePP] (0,0) rectangle (1,1);
            \begin{scope}
                \clip (-1,-1) rectangle (1,1);
                \foreach \r in {0.35,0.7,1.05,1.4} { \draw[PairedE, thin] (\dxo,\dyo) circle (\r); }
            \end{scope}
            \draw[dashed] (-1,0) -- (1,0);
            \draw[dashed] (0,-1) -- (0,1);
            \draw[thick] (-1,-1) rectangle (1,1);
            \node[font=\scriptsize] at (-0.5,-0.5) {$-\,-$};
            \node[font=\scriptsize] at (0.5,-0.5) {$+\,-$};
            \node[font=\scriptsize] at (-0.5,0.5) {$-\,+$};
            \node[font=\scriptsize] at (0.5,0.5) {$+\,+$};
            \fill[PairedF] (\dxo,\dyo) circle (2.2pt);
            \node[anchor=south, font=\small, align=center] at (0,1.05) {\name};
            \node[anchor=north, font=\scriptsize] at (-1,-1) {$-1$};
            \node[anchor=north, font=\scriptsize] at (1,-1) {$1$};
            \node[anchor=north, font=\small] at (0,-1) {$\nu$};
            \node[anchor=east, font=\scriptsize] at (-1,-1) {$-1$};
            \node[anchor=east, font=\scriptsize] at (-1,1) {$1$};
            \node[anchor=east, font=\small] at (-1,0) {$\eta$};
            \draw[->, thick] (0,-1.45) -- (0,-1.95);
        \end{scope}
    }
    \begin{scope}[shift={(-1,-4.4)}]
        \foreach \col in {0,...,3} {
            \begin{scope}[shift={(3.2*\col,0)}]
                \begin{scope}
                    \clip (0,0) rectangle (2.15,2.15);
                    \foreach \r in {0.35,0.7,1.05,1.4,1.75,2.1,2.45,2.8} { \draw[PairedE, thin] (0,0) circle (\r); }
                \end{scope}
                \draw[thick,->] (0,0) -- (2.35,0) node[anchor=west, font=\small]{$u$};
                \draw[thick,->] (0,0) -- (0,2.35) node[anchor=south, font=\small]{$v$};
                \node[anchor=north, font=\scriptsize] at (1,0) {$a$};
                \node[anchor=north, font=\scriptsize] at (2,0) {$2a$};
                \node[anchor=east, font=\scriptsize] at (0,1) {$b$};
                \node[anchor=east, font=\scriptsize] at (0,2) {$2b$};
            \end{scope}
        }
        % coincident: all four pieces fold onto [0,a] x [0,b]
        \draw[thick, fill=elmagGay] (0,0) rectangle (1,1);
        \notechip{0.27}{0.73}{pieceMP}{-+} \notechip{0.73}{0.73}{piecePP}{++}
        \notechip{0.27}{0.27}{pieceMM}{--} \notechip{0.73}{0.27}{piecePM}{+-}
        % edge (1,0): s = + on [0,a], s = - on [a,2a]; both t on [0,b]
        \begin{scope}[shift={(3.2,0)}]
            \draw[thick, fill=elmagGay] (0,0) rectangle (1,1) (1,0) rectangle (2,1);
            \notechip{0.5}{0.73}{piecePP}{++} \notechip{0.5}{0.27}{piecePM}{+-}
            \notechip{1.5}{0.73}{pieceMP}{-+} \notechip{1.5}{0.27}{pieceMM}{--}
        \end{scope}
        % edge (0,1): t = + on [0,b], t = - on [b,2b]; both s on [0,a]
        \begin{scope}[shift={(6.4,0)}]
            \draw[thick, fill=elmagGay] (0,0) rectangle (1,1) (0,1) rectangle (1,2);
            \notechip{0.27}{0.5}{pieceMP}{-+} \notechip{0.73}{0.5}{piecePP}{++}
            \notechip{0.27}{1.5}{pieceMM}{--} \notechip{0.73}{1.5}{piecePM}{+-}
        \end{scope}
        % vertex (1,1): four disjoint rectangles
        \begin{scope}[shift={(9.6,0)}]
            \fill[piecePP] (0,0) rectangle (1,1);
            \fill[pieceMP] (1,0) rectangle (2,1);
            \fill[piecePM] (0,1) rectangle (1,2);
            \fill[pieceMM] (1,1) rectangle (2,2);
            \draw[thick] (0,0) rectangle (2,2) (1,0) -- (1,2) (0,1) -- (2,1);
            \node[font=\scriptsize] at (0.5,0.5) {$+\,+$};
            \node[font=\scriptsize] at (1.5,0.5) {$-\,+$};
            \node[font=\scriptsize] at (0.5,1.5) {$+\,-$};
            \node[font=\scriptsize] at (1.5,1.5) {$-\,-$};
        \end{scope}
        \foreach \col in {0,...,3} { \fill[PairedF] (3.2*\col,0) circle (2.2pt); }
    \end{scope}
\end{tikzpicture}

    \caption{The integration regions of the four touching configurations ($a=b$). Top: the domain $[-1,1]^2$ of
    Theorem~\theoRef{reduction}, split into the pieces $I_s\times I_t$, labelled $s\,t$; the singular point
    $(\nu,\eta)=(d_x,d_y)$ (red) is the centre, the midpoint of an edge or a corner, and the circles are level lines
    of $\rho_{\offs}$. Bottom: the same pieces in the variables $u=\sigma_sa(\nu-d_x)$, $v=\sigma_tb(\eta-d_y)$ of
    Lemma~\lemmaRef{pieces}. In a direction with $d=0$ both pieces fold onto $[0,a]$; with $d=1$ they lie side by side
    on $[0,a]$ and $[a,2a]$. The singular point is the origin, a corner of every piece that contains it.}
    \label{fig:regions}
\end{figure}

\begin{table}[ht]
    \centering
    \caption{The weights~\eqref{eq:weights} on the pieces of Lemma~\lemmaRef{pieces} ($x$-direction).}
    \label{tab:pieces}
    \renewcommand{\arraystretch}{1.35}
    \begin{tabular}{ccccc}
    \toprule
    $d_x$ & $0$ & $0$ & $1$ & $1$ \\
    piece & $I_{-}$, $\sigma=-1$ & $I_{+}$, $\sigma=+1$ & $I_{-}$, $\sigma=-1$ & $I_{+}$, $\sigma=-1$ \\
    $u$ & $[0,\,a]$ & $[0,\,a]$ & $[a,\,2 a]$ & $[0,\,a]$ \\ \midrule
    $W_{00}$ & $\displaystyle \frac{a - u}{a}$ & $\displaystyle \frac{a - u}{a}$ & $\displaystyle \frac{2 a - u}{a}$ & $\displaystyle \frac{u}{a}$ \\[2ex]
    $W_{10}$ & $\displaystyle \frac{\left(a - u\right)^{2}}{2 a^{2}}$ & $\displaystyle \frac{\left(a - u\right) \left(a + u\right)}{2 a^{2}}$ & $\displaystyle \frac{\left(2 a - u\right)^{2}}{2 a^{2}}$ & $\displaystyle \frac{u \left(2 a - u\right)}{2 a^{2}}$ \\[2ex]
    $W_{01}$ & $\displaystyle \frac{\left(a - u\right) \left(a + u\right)}{2 a^{2}}$ & $\displaystyle \frac{\left(a - u\right)^{2}}{2 a^{2}}$ & $\displaystyle \frac{u \left(2 a - u\right)}{2 a^{2}}$ & $\displaystyle \frac{u^{2}}{2 a^{2}}$ \\[2ex]
    $W_{11}$ & $\displaystyle \frac{\left(2 a + u\right) \left(a - u\right)^{2}}{6 a^{3}}$ & $\displaystyle \frac{\left(2 a + u\right) \left(a - u\right)^{2}}{6 a^{3}}$ & $\displaystyle \frac{\left(a + u\right) \left(2 a - u\right)^{2}}{6 a^{3}}$ & $\displaystyle \frac{u^{2} \left(3 a - u\right)}{6 a^{3}}$ \\[2ex]
    \bottomrule
\end{tabular}

\end{table}

\theoTitle{Touching moments}{touching} Let $p\in\{-1,1,3\}$, $\offs=(d_x,d_y)\in\{0,1\}^2$ and $\alpha,\beta\in\{0,1\}$.
On the pieces $s,t\in\{-,+\}$ write, with Table~\ref{tab:pieces},
\begin{equation}
    W_{\alpha\beta}\big|_{I_s} = \sum_{m=0}^{3}c^{s}_{m}\,u^{m}, \qquad
    W_{00}\big|_{I_t} = \sum_{n=0}^{1}e^{t}_{n}\,v^{n},
\end{equation}
and let $A_s\times B_t$ be the image of the piece, with $A_s$ equal to $[0,a]$ or $[a,2a]$ according to $d_x$ and $B_t$
equal to $[0,b]$ or $[b,2b]$ according to $d_y$. Then
\begin{equation}
    \label{eq:touching}
    \Mom{p}{\alpha\beta}{\offs} = \frac{1}{ab}\sum_{s,t\in\{-,+\}}\;\sum_{m=0}^{3}\sum_{n=0}^{1}
    c^{s}_{m}\,e^{t}_{n}\;\cornerdiff_{A_s\times B_t}\Gf{p}{m}{n},
\end{equation}
where, with $G=\Gf{p}{m}{n}$, which vanishes on the axes,
\begin{equation}
    \label{eq:fourpoints}
    \begin{aligned}
        \cornerdiff_{[0,a]\times[0,b]}G   &= G(a,b),\\
        \cornerdiff_{[a,2a]\times[0,b]}G  &= G(2a,b)-G(a,b),\\
        \cornerdiff_{[0,a]\times[b,2b]}G  &= G(a,2b)-G(a,b),\\
        \cornerdiff_{[a,2a]\times[b,2b]}G &= G(2a,2b)-G(2a,b)-G(a,2b)+G(a,b).
    \end{aligned}
\end{equation}
Every touching moment is therefore an exact, finite combination of $\Gf{p}{m}{n}$ at the four points $(a,b)$, $(2a,b)$,
$(a,2b)$, $(2a,2b)$. The $N$-moments follow from Lemma~\lemmaRef{symmetries}(iii) with $a$ and $b$ exchanged.

\vspace{\baselineskip}
\proof
By Theorem~\theoRef{reduction}, $\Mom{p}{\alpha\beta}{\offs}$ is the sum over the four pieces $I_s\times I_t$. On a piece
substitute $u=\sigma_s a(\nu-d_x)$ and $v=\sigma_t b(\eta-d_y)$: then $\D\nu\D\eta=\D u\D v/(ab)$,
$\rho_{\offs}=\sqrt{u^2+v^2}$, the weights become the polynomials of Lemma~\lemmaRef{pieces}, and the piece is
$\int_{A_s}\int_{B_t}\sum_{m,n}c^s_me^t_nu^mv^n(u^2+v^2)^{p/2}\D v\D u$. Apply~\eqref{eq:corner} term by term; the
corners on the axes contribute zero by Proposition~\propRef{gproperties}(i). Each piece is absolutely integrable
($p>-2$), so no limit is involved. \qed

\vspace{\baselineskip}
\noindent\textbf{Remark.} The same corner-sum construction works for every offset, but for distant cells it gives a
sum of large, nearly cancelling terms, whereas \GL{} converges exponentially there. It is meant for the touching
offsets only.

\coroTitle{Homogeneity}{homogeneity} $\Mom{p}{\alpha\beta}{\offs;\lambda a,\lambda b}=\lambda^p\,\Mom{p}{\alpha\beta}{\offs;a,b}$
for $\lambda>0$, so the static moments of a grid depend, up to the factor $a^p$, on the aspect ratio $b/a$ only.

\vspace{\baselineskip}
\proof In~\eqref{eq:reduction}, $\rho_{\offs}$ scales by $\lambda$ and the weights do not depend on $a$ and $b$. \qed

\propTitle{Agreement with the coincident-cell formulas of the thesis}{coincident} Consider the integrals
$\mathcal{I}^{(-1)}$ and $\mathcal{I}^{(1)}$ of the appendix of~\cite{machacek2025method}. Then
\begin{equation}
    \begin{aligned}
    (ab)^2\Mom{-1}{00}{\V{0}} &= \mathcal{I}^{(-1)}[1],\\
    (ab)^2\Mom{-1}{11}{\V{0}} &= \mathcal{I}^{(-1)}[\lambda_1\lambda_1'],\\
    (ab)^2\bigl(\Mom{-1}{10}{\V{0}}-\Mom{-1}{11}{\V{0}}\bigr) &= \mathcal{I}^{(-1)}[\lambda_1(1-\lambda_1')],\\
    (ab)^2\Mom{1}{00}{\V{0}} &= \mathcal{I}^{(1)}[1],\\
    (ab)^2\Mom{1}{11}{\V{0}} &= \mathcal{I}^{(1)}[\lambda_1\lambda_1'].
    \end{aligned}
\end{equation}
The first one reads
\begin{equation}
    \begin{aligned}
    \Mom{-1}{00}{\V{0}} &= \frac{1}{3a^2b^2}\Bigg(2a^3+2b^3-2\sqrt{a^2+b^2}\left(a^2+b^2\right)\\
    &\phantom{{}=\frac{1}{3a^2b^2}\Bigg(}+6a^2b\log\left(\frac{\sqrt{a^2+b^2}+b}{a}\right)
    +6ab^2\log\left(\frac{\sqrt{a^2+b^2}+a}{b}\right)\Bigg).
    \end{aligned}
\end{equation}

\vspace{\baselineskip}
\proof By Corollary~\coroRef{homogeneity} it suffices to take $a=1$ and keep $b>0$ symbolic. The difference of the two
sides, written with~\eqref{eq:touching} and the cited formula, simplifies to zero exactly (SymPy). \qed

\vspace{\baselineskip}
\noindent\textbf{Remark (thesis appendix).} The values of the formulas in the appendix of~\cite{machacek2025method} are
correct, but three labels are misprinted: the integrands of the $\mathcal{I}^{(1)}$ entries are written with $1/R$
instead of $R$; the last $\mathcal{I}^{(1)}$ entry repeats the label $[\lambda_1(1-\lambda_1')]$ instead of
$[\lambda_2(1-\lambda_2')]$; and the integrand of $\mathcal{I}^{(-1)}[\lambda_2(1-\lambda_2')]$ contains $(u+\lambda_2')$
instead of $(v+\lambda_2')$.
